\documentclass[10pt]{article}
\usepackage[margin=20mm]{geometry}
\usepackage{amsmath,amssymb,mathtools,fontspec,unicode-math,xurl,hyperref}
\IfFontExistsTF{DejaVu Serif}{\setmainfont{DejaVu Serif}}{\setmainfont{Latin Modern Roman}}
\setmathfont{latinmodern-math.otf}
\setlength{\parindent}{0pt}\setlength{\parskip}{6pt}\sloppy
\title{Initial scientific review: M25-0096
Counting pseudo-Anosovs as weakly contracting isometries}
\author{Reviewing assistant}
\date{26 October 2026}
\begin{document}
\maketitle
Abstract: One confirmed Major finding concerns the stated scope of Theorem C and Theorem 7.4: the virtually cyclic case is not excluded, and the infinite dihedral group gives an explicit counterexample satisfying their hypotheses. The later Theorem 7.9 does contain the necessary exclusion. The mapping-class-group argument was traced through weak contraction, alignment, concatenation, and the insertion-fibre count. The source constant in Proposition 4.1 has been checked and a prior numerical criticism withdrawn. The remaining local arithmetic clarifications and an unresolved verification of the arbitrary-number-of-insertions step are separated from the confirmed finding. The report does not refute the mapping-class-group theorem or the correctly restricted extensions. Human scientific approval remains pending.\par
\section*{Bound version and inspected scope}
This review concerns Inhyeok Choi’s arXiv:2408.00603v3, submitted 12 October 2025; the bound PDF is dated 14 October 2025 and has 38 pages and 646,186 bytes. The journal identity is DOI 10.1007/s00222-025-01358-5, Inventiones mathematicae 242 (2025), 337–386, published online 22 August 2025. The journal Version of Record was not acquired. Neither metadata linkage nor the repository’s description of v3 as accepted establishes equality with the journal text.\par
All locations below are author-version PDF page numbers. Full extracted text on pp. 1–36 was read. Pages 1–3 and 7–36 (33 pages) were additionally actually viewed as original page images, including every main theorem, formal proof section, and Appendix A. Pages 4–6 were text-reviewed but not visually inspected; pp. 37–38 were not independently reviewed. The full 170,741-byte author TeX was parsed statically, with targeted mathematical source contexts compared to the PDF. Parsing an index or source inventory is not scientific review. No author source was compiled or executed.\par
The eight manifest files were verified against their recorded sizes and SHA-256 values. Author TeX, bibliography and archive are typesetting material, not released computational code. Absence of released code is not a defect. An independently written normal-form calculation was actually executed and retained; the general counterexample proof below does not depend on finite numerical sampling.\par
\section*{Claims, hypotheses and dependency audit}
Theorem A (p. 2) and Theorem 5.3 (p. 18) fix an arbitrary finite generating set of the mapping class group of the specified finite-type hyperbolic surface, excluding the three-punctured sphere. The stronger statement controls the union of the small curve-complex translation-length event and the event \(\tau_S(g)\le0.33R\). Constants may depend on the generating set. They are not uniform over all generating sets. Word-ball counting is not interchangeable with random-walk counting or counting ambient-metric balls.\par
\subsection*{Projection and concatenation, pp. 7–18 and 34–36}
The fixed pseudo-Anosov has a geodesic axis in the curve complex, with \(D_C>0\), and the orbit map is Lipschitz. Fact 2.2 uses two different inputs: non-elementarity supplies bounded connecting letters, while WPD turns long overlapping oriented axis segments into finitely many possibilities in the word metric. This distinction matters in Finding F1: WPD alone does not supply non-elementarity.\par
Lemma 2.4 follows by choosing an orbit point near the projected point and using the comparison \(d_S(a,b)\le(D_S/D_C)d_C(ax_0,bx_0)\) along the chosen orbit. Proposition 2.3 is an external weak-contraction input; its original subsurface-projection proof was not reconstructed here. The alignment convention in Definition 3.2 controls terminal and initial projections in the prescribed orientation. Facts 3.3–3.5, Lemmas 3.6–3.7 and the Appendix A projection estimates were traced, including the requirement that interior segments exceed the alignment-error thresholds.\par
Proposition 4.1 chooses a crossing component of a word-geodesic whose curve-complex projections lie near the axis midpoint. If the component stays far away in the word metric, weak contraction forces many long steps before the projected crossing can finish. Its choice \(T=103K_{\mathrm{map}}^3/D_C\) satisfies the displayed comparison; the previously reported numerical adjustment is withdrawn below. Proposition 4.3 then distributes errors with weights \(30^{-j}\). The induction’s left and right coefficient gaps dominate the extra contribution from the newly chosen first point. In Corollary 4.5, \[2\sum_{j\ge1}30^{-j}=\frac2{29}<\frac1{14};\] the ordered points on the word-geodesic therefore give the required bound on the sum of distances to the orbit midpoints.\par
Corollary 4.6 combines linear progress of the aligned segments with the fact that the two relevant word-geodesic pieces concatenate to a geodesic. For the range used in the fibre argument, summing \(D_i+D_i'\ge 2K_{\mathrm{map}}(i-L')\) gives a quadratic lower bound. The displayed sum proves this directly when the number of segments is at least \(L'\). For smaller numbers, the fibre estimate is already bounded by a constant; one should use the positive part \((n-L')_+^2\) in a universally stated lower bound. This audit does not claim that the printed stronger small-\(n\) bound has been independently established.\par
\subsection*{Counting and translation-length bounds, pp. 18–24}
Fekete’s lemma applies to the submultiplicative word-ball counts, and positive exponential growth is essential for the inner-ball ratio. For a fixed positive growth exponent \(\lambda\), choose the upper growth-rate error sufficiently smaller than \((1-0.99)\lambda\); then \(|B_S(0.99n)|/|B_S(n)|\) decays exponentially.\par
In Lemma 5.5, an aligned orbit near one quarter of the word length either iterates to produce positive translation length, or yields a conjugacy expression with conjugator length at most \(0.31n\) and conjugated element length at most \(0.57n\). The conjugator is counted once, not twice: it also determines its inverse. Thus \(0.31+0.57<1\) makes the latter set exponentially negligible. The retained exact arithmetic check gives \[0.86(0.99-0.6-0.001)-0.334=\frac{27}{50000}>0,\qquad
0.334\cdot0.99-0.33=\frac{33}{50000}>0.\] These margins remove a superficial concern about the non-strict inequalities at the translation-length cutoff. The lower distance in the definition of \(\mathcal A_{\mathrm{thick}}\) is \(0.25\|g\|_S\), not necessarily \(0.25n\); the later estimates can and do retain the former bound.\par
The insertion map \(F_n\) removes a fixed-length block and replaces it by \(s\phi^{L_{\mathrm{map}}}t\), so its image remains in \(B_S(n)\). Overlapping trimmed intervals on \([x_0,Ux_0]\) are clustered by WPD, with a number of word choices independent of \(n\). A maximal disjoint family is aligned. If its last segment aligns with the translated endpoint, the original input lies in \(\mathcal A_{\mathrm{thick}}\), contrary to the domain definition. Otherwise Corollary 4.6 bounds the disjoint family by \(O(\sqrt n)\). There are \(\Theta(n)\) insertion positions, so the double count gives an \(O(n^{-1/2})\) exceptional proportion. The fixed constants and the finitely many small radii can be absorbed into the theorem’s \(K\).\par
\subsection*{Two insertions and extensions}
Lemmas 6.1–6.2 count a last-stage prefix/block/suffix triple and then its possible predecessors, each with \(O(\sqrt n)\) possibilities. The midpoint case split in Lemma 6.2 uses the original geodesic word and the original input’s exclusion from \(\mathcal A_{\mathrm{thick}}\); an arbitrary intermediate word is not automatically another valid input. This is why a complete multi-stage induction needs to retain those data. The two-insertion argument gives an \(O(n^{-1})\) bound after correcting a harmless counting coefficient below. The proof of Theorem 6.3 explicitly writes out only \(k=2\); verification of the all-\(k\) step is recorded separately as U1.\par
The HHG and quasi-isometry extensions require the external results on maximal hyperbolic spaces, Morse detection, and weak contraction. Lemma 7.8 transports projections through quasi-isometries; the constants must be enlarged together so that the additive errors are absorbed. For the CAT(0) extension, geometric action gives word/ambient metric comparison and properness gives WPD; the rank-one/strong-contraction equivalence is an external input. These external theorems were not re-proved.\par
For braids, the argument first counts cosets of the center. If \(A=\max_{s\in S}|\rho(s)|\) for exponent sum \(\rho\), then \(|\rho(g)|\le A\|g\|_S\), and \(\rho(\Delta^2)\ne0\) bounds each central coset’s intersection with \(B_S(R)\) by \(O(R)\). Quotient word balls are images of the original word balls. Thus a quotient exceptional bound \(O(R^{-q})\) yields \(O(R^{1-q})\) upstairs, and an arbitrary polynomial exponent is needed to deduce the final superpolynomial statement. The literal phrase that word norm only decreases under homomorphisms must be read with the appropriate generating-set Lipschitz constant. No claim that a braid itself is a Morse element is required.\par
\section*{F1: confirmed Major — missing non-virtually-cyclic hypothesis}
\subsection*{Exact location and claim}
Theorem C, p. 3 (author TeX lines 167–172), allows every group quasi-isometric to a well-behaved HHG containing Morse elements. Theorem 7.4, pp. 28–29 (TeX lines 1471–1476), assumes a finitely generated group with a WPD loxodromic isometry weakly contracting by means of the hyperbolic space. Neither statement excludes virtually cyclic groups. Theorem C claims an \(O(R^{-1/2})\) non-Morse proportion; Theorem 7.4 claims superpolynomial genericity of loxodromics.\par
The displayed numerator of Theorem 7.4 counts loxodromics instead of non-loxodromics. That sign-of-event error is not the basis of this finding: the counterexample below defeats the intended genericity conclusion even after the numerator is corrected.\par
\subsection*{Full counterexample and every hypothesis}
Let \[G=D_\infty=\langle r,s\mid s^2=1,\ srs=r^{-1}\rangle,\qquad
S=\{r,s\},\qquad X=\mathbb R.\] Use the action \(r(x)=x+1\), \(s(x)=-x\) and basepoint \(0\). The space is geodesic and \(0\)-hyperbolic, and the action is isometric. Every element has a unique normal form \(r^m s^\epsilon\), \(m\in\mathbb Z\), \(\epsilon\in\{0,1\}\). The orbit of \(\phi=r\) is the integer line, hence a geodesic orbit and loxodromic.\par
For WPD, the condition \(d_X(0,h0)<L\) already forces \(|m|<L\) and permits only two values of \(\epsilon\). It gives a finite bound independent of the second condition, so any positive \(N\), with sufficiently large \(M\), works in Definition 7.3. This verifies WPD directly, rather than inferring non-elementarity from it.\par
For weak contraction as Definition 7.2 states it, put \[\pi_\gamma(r^m s^\epsilon)=r^m,\qquad \gamma=\{r^j:j\in\mathbb Z\}.\] This projection fixes \(\gamma\), and its image under the orbit map is exactly the nearest projection of \(h0=m\) to the integer orbit. It is therefore a valid projection by means of \(X\). Also \[d_S(r^m s^\epsilon,\gamma)=\epsilon\le1.\] Taking the weak-contraction threshold \(K'=2\) makes the required condition for points farther than \(K'\) vacuous. The orbit is nevertheless an actual geodesic. The theorem has no requirement that the complement of the orbit have unbounded distance, so this is a legitimate instance of all its assumptions.\par
The word length is exactly \[\|r^m s^\epsilon\|_S=|m|+\epsilon.\] The upper bound is its normal form. For the lower bound, each \(r^{\pm1}\) changes the integer coordinate by at most one, and a word representing \(\epsilon=1\) must contain at least one \(s\). For every integer \(n\ge1\) there are \(2n+1\) translations with \(|m|\le n\) and \(2n-1\) reflections with \(|m|\le n-1\). The nonzero translations are loxodromic; every reflection squares to the identity, and the identity is also non-loxodromic. Hence \[|B_S(n)|=4n,\qquad
|\{g\in B_S(n):g\text{ is non-loxodromic}\}|=2n,\qquad
\frac{2n}{4n}=\frac12.\] The ratio cannot be at most \(K n^{-k/2}\) for any fixed \(K\) and any \(k>0\) at all large \(n\). The literal loxodromic numerator also has ratio \(1/2\). Actual independent enumeration of balls of radii 1–20 agrees with these formulas; the formulas above prove the assertion for every \(n\).\par
For Theorem C, choose \(H=\mathbb Z\) with its single-domain HHG structure: the top hyperbolic space is its Cayley line and the projection is the identity. This is normalized; bounded-domain dichotomy holds; unbounded products imposes no condition because there is no proper domain. The top space has a cobounded isometry group and uniform bi-infinite geodesics, so it is well-behaved in the cited Definition 6.9. The definition and the proper-domain restriction were checked in the 2023 primary preprint of Goldsborough–Hagen–Petyt–Sisto, pp. 20, 23 (with Russell’s appendix). No equality with its later journal version is assumed.\par
The generator of \(\mathbb Z\) is Morse, since its orbit is the whole Cayley line and quasi-geodesics in a line stay within a bound depending only on their quasi-geodesic constants. Finally, \[f:D_\infty\longrightarrow\mathbb Z,\qquad f(r^m s^\epsilon)=m\] is surjective and changes distances by at most one: \(d_S(r^m s^\epsilon,r^l s^\eta)=|m-l|+(\epsilon\mathbin{\mathrm{xor}}\eta)\). Thus \(G\) is quasi-isometric to this well-behaved \(H\). Reflections and the identity have bounded cyclic orbits and are not Morse elements. Nonzero translations have geodesic, coarsely dense orbits in the quasi-line, and are Morse. The non-Morse proportion is again exactly \(1/2\), contradicting Theorem C’s bound.\par
\subsection*{Counterevidence, impact and repair limits}
Theorem 7.6 (p. 29), Theorem 7.9 (p. 31; TeX line 1541), Theorem 7.13 (p. 32) and Theorem 7.15 (p. 33) explicitly require the relevant group to be non-virtually cyclic. This is strong internal evidence of the intended restriction. The non-elementarity argument in Fact 2.2(1), p. 8, and positive growth exponent in Lemma 5.1, p. 18, also identify exactly why the proof cannot be transferred to this example.\par
Add the non-virtually-cyclic hypothesis to Theorem C and Theorem 7.4, and correct Theorem 7.4’s exceptional event. For Theorem C this makes its scope agree with Theorem 7.9. This is a restriction of an advertised theorem, not a proof of its unrestricted form. The counterexample does not address the mapping-class-group result, the non-virtually-cyclic CAT(0) cases, or the correctly restricted Theorem 7.9. It also does not independently verify the general multi-insertion step or the external HHG machinery. Severity is Major because a principal quantitative genericity statement is false on an allowed class, even though the scope repair is short.\par
\section*{Reconstructed estimates and unresolved verification}
\subsection*{Local numerical corrections, without an additional core-failure claim}
In Proposition 4.1, p. 12, write \(A=K_{\mathrm{map}}\). The source chooses \(T=103A^3/D_C\), as confirmed directly in both the original PDF page image and author TeX line 596. The comparison is valid: \[\begin{aligned}
(10^6-1)A^6-6000A^3(1+D_CT)
&=A^3(381999A^3-6000)>0
\end{aligned}\] for \(A\ge1\); at \(A=1\) the margin is \(375999\). Also \[D_CT-A=103A^3-A\ge102A^3.\] No change of \(T\) is required. The previous numerical criticism and proposed \(T=150A^3/D_C\) repair resulted from reading \(103\) as \(1000\) and are withdrawn. This correction does not affect Finding F1. Discrete word-geodesic steps can be rounded down while retaining the original projected-traversal slack \(103A^3-A\) before weakening the bound; this is separate bookkeeping, not a defect in the source’s choice of \(T\).\par
In (5.4), p. 23, Proposition 4.1 gives \(0.01(d_S(\mathrm{id},q)+d_S(q,p))\) and the final lower bound retains a factor \(0.98\) on both terms. The useful conclusion survives: \[d_S(\gamma_1,\gamma_M)
\le\left(\frac{0.28}{0.98}-0.27\right)n
=\frac{11}{700}n<0.016n.\] For the subsequent lower distance to the translated endpoint, \(0.274n-0.016n-O(1)=0.258n-O(1)\) suffices to exceed \(0.25\|g_1\|_S\) at large \(n\); the displayed stronger \(0.26n\) is unnecessary. Distances and diameters of these fixed-length orbit segments differ by \(O(1)\), which must be retained before taking large \(n\).\par
In the two-cut index set on p. 24, let \(m\) be the number of available consecutive integers and let \(D=2K_{\mathrm{map}}L_{\mathrm{map}}+3\). The actual number of pairs with \(j-i>D\) is \[\frac{(m-D-1)(m-D)}2 \quad (m>D+1).\] Thus the coefficient-one product printed in the last display on p. 27 is not the literal count. It remains \(\Theta(n^2)\) for fixed \(D\), which is all the stated \(O(n^{-1})\) argument needs. These repairs concern local bookkeeping; they are not additional counterexamples to the main genericity theorem.\par
\subsection*{U1: Needs Evidence — the arbitrary-number-of-insertions step}
Locations: the opening of Section 6, p. 24; Theorem 6.3 and its proof, p. 27; transfer to Corollary 7.20, p. 34. The proof explicitly treats \(k=2\) and states that the same idea gives the general case. The natural desired extension has \(\Theta(n^k)\) separated choices and an \(O_k(n^{k/2})\) fibre bound. To certify it one needs intermediate records retaining the original geodesic, removed blocks, fixed suffixes and the original exclusion from \(\mathcal A_{\mathrm{thick}}\), with the analogue of both cases of Lemma 6.2 at every stage.\par
The two-insertion construction provides substantial positive evidence for such an induction, and no all-hypotheses counterexample to the superpolynomial assertion was found. This review did not complete that induction with all alignments and endpoint constants. Accordingly U1 is an unresolved verification item, not a confirmed Major or Minor defect and not evidence that the superpolynomial theorem is false. Its relevance is substantive: the \(O(R)\) central-fibre loss in the braid argument requires an exponent beyond the explicitly written \(k=2\) case.\par
\section*{Supporting checks}
The actual calculations used reviewer-written Python normal forms and exact rational arithmetic, with exit code zero; no author code was run. The complete author TeX static record reports 1,793 lines, 40 parsed theorem-like environments and 92 labels, its package/bibliography commands, and targeted theorem contexts. Those counts only describe the inspection input. The source images and page-specific observations are retained separately.\par
The repository record independently confirms v3 and states that it includes corrections based on a referee report. A limited title/DOI correction search did not identify a separate article-specific erratum; this is not proof that no correction exists. The external definition check used https://arxiv.org/abs/2309.07013, v1, via the University of Bristol primary preprint copy, Definition 6.9 and the definitions at the start of Section 6. Only those definitions were checked for F1; the full external paper was not audited.\par
The final report is compiled from the current TeX using the authorized local Tectonic 0.17.0. Source and the complete cached dependency set must have identical pre/post hashes; the produced PDF, compiler, logs and inspected page images are hash-bound by provenance and delivery state. Page rendering alone does not count as inspection.\par
This initial review has been corrected during independent focused acceptance, with the predecessor retained. Finding F1 and the explicit limitation U1 remain as stated. Human scientific approval remains pending. `human\_reviews=[]`; `scientific\_decision=pending`.\par
\end{document}
